% -*-LaTex-*-
% Web CV for Matthew J. Holman (from the 4 October 2026 nomination CV, phone number removed).
% Preamble deliberately identical to holman-bio-sketch-26-2p_edited.tex, which
% is known to compile.  Only geometry, array and tabularx are used.
\documentclass[11pt,letterpaper]{article}

\usepackage[margin=1.0in]{geometry}
\usepackage{array}
\usepackage{tabularx}

\newcommand{\sect}[1]{\vspace{0.14in}\noindent{\bf #1}\vspace{0.06in}\\}

\pagestyle{plain}

\begin{document}

\begin{center}
{\large {\bf MATTHEW J. HOLMAN}}
\end{center}

\noindent
{\renewcommand{\arraystretch}{1.05}%
\begin{tabularx}{\textwidth}{@{}>{\raggedright\arraybackslash}p{3.9in}>{\raggedleft\arraybackslash}X@{}}
Center for Astrophysics $\vert$ Harvard \& Smithsonian & {\it mholman@cfa.harvard.edu}\\
Mail Stop \#51 & ORCID: 0000-0002-1139-4880\\
60 Garden Street & \\
Cambridge, MA 02138 &
\end{tabularx}}

\sect{Education and Training}
\noindent
{\small
\renewcommand{\arraystretch}{1.08}%
\begin{tabularx}{\textwidth}{@{}>{\raggedright\arraybackslash}p{3.25in}>{\raggedright\arraybackslash}X>{\raggedleft\arraybackslash}p{1.25in}@{}}
Massachusetts Institute of Technology & Mathematics & SB, 1989\\
Massachusetts Institute of Technology & Planetary Sciences & PhD, 1994\\
National Astronomical Observatory of Japan & Solar System Dynamics & Postdoc, 1994--1995\\
Canadian Institute for Theoretical Astrophysics & Solar System Dynamics & Postdoc, 1995--1997\\
Harvard University & Comput. Sci. \& Eng. & SM, 2018
\end{tabularx}}

\sect{Appointments}
\noindent
{\small
\renewcommand{\arraystretch}{1.08}%
\begin{tabularx}{\textwidth}{@{}>{\raggedright\arraybackslash}p{2.9in}>{\raggedleft\arraybackslash}X@{}}
Smithsonian Astrophysical Observatory & {\it Senior Astrophysicist.} 2011--present\\
Harvard University & {\it Lecturer on Astronomy.} \\
Minor Planet Center & {\it Director.} 2015--2021\\
Smithsonian Astrophysical Observatory & {\it Associate Director.} 2004--2014\\
Smithsonian Astrophysical Observatory & {\it Astrophysicist.} 1997--2011
\end{tabularx}}

\sect{Honors and Awards}
\noindent
{\small
AAAS Newcomb Cleveland Prize, for the best paper published in {\it Science}, 2000\\
Asteroid 3666 Holman, 1999\\
Boris Garfinkel Lectureship, Yale University, 2018\\
Bok Center Excellence in Teaching Award, Harvard University, 2017\\
NCU--Delta Young Astronomer Lectureship Award, 2012\\
Cover article of {\it Science}, 1 October 2010\\
NASA Group Achievement Award, Kepler Launch and Commissioning Team, 2010\\
NASA Group Achievement Award, EPOXI Project Team, 2009\\
Editor's Letter of Commendation for Excellent Reviewing, {\it Icarus}, 1998\\
Japan Society for the Promotion of Science Postdoctoral Fellowship, 1994\\
National Research Council Research Associateship (declined), 1994\\
National Science Foundation Graduate Fellowship, 1990--1993
}

\sect{Principal Research Contributions}
\noindent
{\small
\begin{itemize}\itemsep2pt \parskip0pt
\item {\it $N$-body integrators.} With Jack Wisdom, introduced the symplectic
$n$-body mapping now known as the Wisdom--Holman map, an order of magnitude
faster than prior methods and since adopted as the standard technique for
long-term integration of the solar system.  Led the
development of ASSIST, an ephemeris-quality test-particle integrator built for
the precision that modern surveys require.

\item {\it Transit timing variations.} With Norman Murray, proposed the use of
deviations in transit times to detect planets far below the mass reachable from
transit depth alone, and in 2010 led the discovery of Kepler-9, the first
planetary system confirmed by that method.  Transit timing is now a standard
technique in exoplanet science.

\item {\it Faint trans-Neptunian discovery from space.} The small end of the
trans-Neptunian size distribution records the collisional history of the outer
solar system.  The bodies that carry that record are too faint for
ground-based surveys.  Beginning with the {\it Hubble Space Telescope} survey
that first measured the small end and found fewer small bodies than collisional
models predicted (Bernstein et al.\ 2004).  Has pursued the same question through
deeper searches of those {\it HST} data and a combined {\it JWST} and {\it HST} program that in 2026
reached diameters near 10~km.  Yet deeper searches with {\it JWST} are underway.  Will use the {\it Nancy Grace Roman Space Telescope}, through a funded archival search (PI Napier) of the
Galactic Bulge Time Domain Survey, to discover $\sim1000$ faint TNOs.

\item {\it Stability and the origin of chaos.} Mapped the stable and unstable
regions among the giant planets, together with the dynamical pathway that
supplies short-period comets, and established where planets can persist in
binary systems --- a criterion published before any such planet was known and
later applied to the circumbinary planets found by Kepler.  Then, with Norman
Murray, identified the dynamical mechanism responsible for chaos in asteroid and
planetary orbits --- the overlap of high-order and three-body mean-motion
resonances --- in a series of papers running from the outer asteroid belt
({\it The Astronomical Journal}, 1996, 1997, 1998) to the outer solar system
({\it Science}, 1999) and the role of chaotic resonances generally ({\it
Nature}, 2001).%  Invited, with Lecar and Franklin, to write the {\it Annual
%Review of Astronomy and Astrophysics} treatment of the subject (2001).

\item {\it Discovery at survey scale.} Developed HelioLinC, the linking
algorithm that makes minor-planet discovery possible in sparse survey data
without nightly pairs.  A version of it, developed by Ari Heinze, now runs in the
Vera C.\ Rubin Observatory's solar system pipeline, where it is expected to account for the majority of the
survey's discoveries over its ten-year planned operations.
\end{itemize}
}

\sect{Discoveries}
\noindent
{\small
Credited with the discovery or co-discovery of one moon of Jupiter, one of
Saturn, seven of Uranus, and five of Neptune.  In the past year, discovered
approximately 1{,}500 trans-Neptunian objects in Rubin Observatory data, all
reported to the Minor Planet Center, against roughly 6{,}000 known previously --
a one-quarter increase in the known population of the outer solar system.
}

\sect{Publications}
\noindent
{\small
244 refereed publications.  46{,}995 citations and $h = 91$ (NASA ADS, 4 October
2026).  Twelve publications selected as most significant are listed separately.
}

\sect{Training and Mentoring}
\noindent
{\small
Mentored 14 postdoctoral researchers, 8 doctoral students, 3 master's
students, and 28 undergraduate researchers, and served on roughly 15 doctoral
supervisory committees.

Former postdoctoral mentees include Konstantin Batygin (Caltech), B.~Scott
Gaudi (Ohio State), Joshua Winn (Princeton), Daniel Fabrycky (Chicago), Eric
Ford (Penn State), Ruth Murray-Clay (UC Santa Cruz), Darin Ragozzine (Brigham
Young), and Matthew Payne (Director, Minor Planet Center).  Former undergraduate
researchers include Shannon Curry (LASP, Colorado), Nikole Lewis (Cornell),
Carolyn Ernst (JHU Applied Physics Laboratory), and Wesley Fraser (NRC
Herzberg).
}

\sect{Selected Service}
\noindent
{\small
Users Committee for NSF's NOIRLab, 2024--present\\
Rubin Observatory Users Committee: Chair 2022--2024, Member 2024--2025\\
Rubin Observatory Solar System Commissioning Team, 2021--present\\
%LSST Solar System Science Collaboration, Commissioning Team\\
Roman Early Definition Survey Assessment Committee, 2023--2024\\
American Association for the Advancement of Science, Astronomy Section,
Member-at-large, 2023--2024\\
Director, Minor Planet Center, 2015--2021\\
Pan-STARRS-1 Science Consortium Board, 2008--2016\\
Kepler Science Council, elected member, 2009--2011\\
Division on Dynamical Astronomy, Brouwer Award Selection Committee, 2001--2003
(Chair, 2003)\\
Co-chair, CfA Colloquium Committee, 2026--present\\
Smithsonian Astrophysical Observatory, Professional Accomplishments Evaluation
Committee, 2026\\
Referee for {\it Nature}, {\it Science}, {\it The Astronomical Journal}, {\it
The Astrophysical Journal}, {\it Icarus}, {\it The Planetary Science Journal},
and {\it Celestial Mechanics and Dynamical Astronomy}; NASA and NSF grant review
panels
}

\end{document}

\sect{Current Research Support}
\noindent
{\small
NASA Roman Large Cycle 1, 2026--, ``An Archival Deep Drilling Kuiper Belt Search
in the Galactic Bulge Time Domain Survey'' (PI Napier, Co-I Holman), \$1.2M
total, \$800{,}000 to SAO.\\[2pt]
NSF AST-2206194, 2022--2027, ``Collaborative Research: A Gravitational Inventory
of the Solar System Using High Precision Minor Planet Astrometry'' (PI Holman),
\$293{,}281.\\[2pt]
Two JWST Cycle 6 proposals submitted as Principal Investigator, September 2026.
}


